% Part: first-order-logic
% Chapter: proof-systems
% Section: tableaux

\documentclass[../../../include/open-logic-section]{subfiles}

\begin{document}

\iftag{FOL}
      {\olfileid{fol}{prf}{tab}}
      {\olfileid{pl}{prf}{tab}}

\olsection{\usetoken{P}{tableau}}

Veel stelsels van afleidingsregels werken met !!{sentence}s die op een
bepaalde manier bij elkaar staan. !!{tableau}s werken met
!!{signed formula}s. !!^a{signed formula} is een paar: een teken voor
een waarheidswaarde ($\True$ of $\False$) en !!a{sentence}:
\[
\sFmla{\True}{!A} \text{ of } \sFmla{\False}{!A}.
\]
!!^a{tableau} is een boom van !!{signed formula}s die naar beneden
vertakt. Bovenaan staan enkele \emph{aannames}. Daaronder komen
!!{signed formula}s die met een afleidingsregel uit een erboven staande
!!{signed formula} volgen. Met elke regel voegen we aan het einde van
een tak een of meer !!{signed formula}s toe. Soms voegen we twee
!!{signed formula}s naast elkaar toe. Dan splitst de tak zich in
tweeën en staat aan elk nieuw uiteinde een van de twee toegevoegde
!!{signed formula}s.

Als we een regel op een samengestelde !!{signed formula} toepassen,
voegen we de directe deel-!!{formula}s ervan toe. Voor elk van de twee
tekens is er een eigen regel. Zo passen we $\TRule{\True}{\land}$ toe
op $\sFmla{\True}{!A \land !B}$. We mogen dan zowel
$\sFmla{\True}{!A}$ als~$\sFmla{\True}{!B}$ toevoegen aan het einde
van elke tak die $\sFmla{\True}{!A \land !B}$ bevat. Met
$\TRule{\False}{!A \land !B}$ mogen we een tak in tweeën splitsen. We
voegen dan $\sFmla{\False}{!A}$ en $\sFmla{\False}{!B}$ naast elkaar
toe. !!^a{tableau} is gesloten als elke tak een paar !!{signed formula}s
met dezelfde formule maar verschillende tekens bevat:
$\sFmla{\True}{!A}$ en $\sFmla{\False}{!A}$.

Voor !!{tableau}s definiëren we de relatie $\Proves$ als volgt.
$\Gamma \Proves !A$ geldt precies als er een eindige verzameling
$\Gamma_0 = \{!B_1, \dots, !B_n\} \subseteq \Gamma$ bestaat waarvoor
er een gesloten !!{tableau} is met de volgende aannames:
\[
\{\sFmla{\False}{!A}, \sFmla{\True}{!B_1}, \dots, \sFmla{\True}{!B_n}\} 
\]
De volgende gesloten !!{tableau} laat bijvoorbeeld zien dat $\Proves
(!A \land !B) \lif !A$:
\begin{oltableau}
  [\sFmla{\False}{(\formula{A} \land \formula{B}) \lif \formula{A}}, just = \TAss
    [\sFmla{\True}{\formula{A} \land \formula{B}}, just = {\TRule{\False}{\lif}[1]}
      [\sFmla{\False}{\formula{A}}, just = {\TRule{\False}{\lif}[1]}
        [\sFmla{\True}{\formula{A}}, just = {\TRule{\True}{\lif}[2]}
          [\sFmla{\True}{\formula{B}}, just = {\TRule{\True}{\lif}[2]}, close]
        ]
      ]
    ]
  ]
\end{oltableau}

Een verzameling $\Gamma$ heet inconsistent in de !!{tableau}calculus als
we een gesloten !!{tableau} kunnen maken met de aannames
\[
\{\sFmla{\True}{!B_1}, \dots, \sFmla{\True}{!B_n}\} 
\]
Daarbij moet $!B_i \in \Gamma$ gelden voor elke gebruikte index.

Evert Beth en Jaakko Hintikka vonden !!{tableau}s in de jaren vijftig
van de twintigste eeuw onafhankelijk van elkaar uit. Raymond Smullyan
maakte de methode daarna eenvoudiger en bekender. Een !!{tableau} is
heel gemakkelijk te gebruiken omdat je bij het opbouwen een vaste reeks
stappen volgt. Daarom kan ook een computer de methode goed uitvoeren.
!!^a{tableau} is vaak wel moeilijk te lezen. Ook is
niet altijd duidelijk hoe het met een bewijs samenhangt. We kunnen de
methode ook voor veel verschillende logica's gebruiken. Voor die
logica's bestaan !!{tableau}systemen. Met !!{tableau}s kunnen we ook
!!{structure}s
vinden die gegeven (verzamelingen) !!{sentence}s vervullen. Als de
verzameling vervulbaar is, kan er geen gesloten !!{tableau} voor zijn.
Elk !!{tableau} heeft dan een open tak. Uit zo'n open tak kunnen we de
vervullende !!{structure} ``aflezen''. Daarvoor moeten we op die tak
wel elke regel hebben toegepast die erop kon worden toegepast. Er is
ook een heel nauw verband met de sequentencalculus. Een gesloten
!!{tableau} is in wezen een !!{derivation} in de sequentencalculus
waarin de stappen dichter op elkaar staan en die ondersteboven is
geschreven.
\end{document}
